\documentclass[11pt,a4paper]{article}
\usepackage[utf8]{inputenc}
\usepackage[spanish]{babel}
\usepackage{amsmath, amssymb}
\usepackage{geometry}
\geometry{top=2.5cm, bottom=2.5cm, left=2.5cm, right=2.5cm}

\begin{document}

\section*{Piense, dialogue y comparta}

\begin{itemize}
    \item[\textbf{1.}] Usted es un jugador de futbol en el equipo de su escuela. En la práctica, patea un balón de \(0.400\,\text{kg}\) del suelo en la yarda 20 de un campo de futbol que está exactamente a lo largo de una recta norte-sur. El vector de velocidad inicial de la pelota está dirigido a \(30.0^\circ\) arriba de la horizontal, hacia el norte, y tiene una magnitud de \(27.0\,\text{m/s}\). Usted ve la pelota cuando se levanta en el aire. Cuando la bola alcanza su punto más alto en su trayectoria parabólica, un águila de \(4.18\,\text{kg}\) volando hacia el sur y a lo largo de una recta horizontal choca con la bola. Suponga que la bola muy inflada hace una colisión elástica con el pico duro del pájaro, y que la pelota rebota por la colisión con un vector de velocidad que es horizontal y hacia el sur. La bola aterriza de nuevo en el punto exacto en el suelo en que fue pateada. 
    \begin{enumerate}
        \item ¿Qué tan rápido volaba el águila? (desprecie la resistencia del aire)
        \item ¿Qué tan rápido y en qué dirección se mueve el águila justo después de la colisión? (suponga que aún no ha ocurrido el aleteo)
        \item En realidad, la colisión no será elástica: algo de la energía cinética se convertirá en otra forma de energía. Suponiendo que algo de la energía cinética se pierde en una colisión inelástica, ¿la velocidad requerida del águila será grande o más pequeña que en el inciso (a)?
    \end{enumerate}
    \vspace*{9cm}

    \item[\textbf{2.}] Dibuje cuidadosamente un triángulo rectángulo sobre un trozo de cartón, tal que uno de sus catetos sea \(30\text{–}40\,\text{cm}\) de longitud y el otro cateto es mucho más corto. Mida el punto medio exacto de cada uno de los tres lados del triángulo, y marque estos tres puntos. Trace una recta a través del triángulo, desde un vértice del triángulo al punto medio de la parte opuesta. Repita para las otras dos esquinas. Las tres rectas se intersecan en el centro de masa del triángulo. Dibuje una cuarta recta perpendicular al cateto más largo, que pase a través del centro de masa, y termine cuando cruza la hipotenusa del triángulo. Perfore un agujero en el cartón justo dentro del borde del triángulo donde la cuarta recta cruza la hipotenusa. Corte con cuidado el triángulo de la cartulina. Ate una cuerda a través del agujero y cuelgue el triángulo de la cuerda. El lado más largo del triángulo debe estar paralelo a la mesa. ¿Por qué debe ser cierto esto? Ahora mida la distancia a lo largo del cateto desde el ángulo más pequeño a la cuarta recta. ¿Qué fracción de todo el cateto más largo es esta distancia?
    \vspace*{7cm}
\end{itemize}

\newpage

\section*{Problemas}

\subsection*{SECCIÓN 9.1 Cantidad de movimiento lineal y SECCIÓN 9.2 Análisis de modelos: Sistema aislado (cantidad de movimiento)}

\begin{itemize}
    \item[\textbf{1.}] Una partícula de masa \(m\) se mueve con cantidad de movimiento de magnitud \(p\).
    \begin{enumerate}
        \item Demuestre que la energía cinética de la partícula es
        \[
        K = \frac{p^2}{2m}.
        \]
        \item Exprese la magnitud de la cantidad de movimiento de la partícula en términos de su energía cinética y masa.
    \end{enumerate}
    \vspace*{6cm}

    \item[\textbf{2.}] Una partícula de \(3.00\,\text{kg}\) tiene una velocidad de
    \[
    (3.00\,\hat{i} - 4.00\,\hat{j})\,\text{m/s}.
    \]
    \begin{enumerate}
        \item Encuentre las componentes \(x\) y \(y\) de su cantidad de movimiento.
        \item Encuentre la magnitud y dirección de su cantidad de movimiento.
    \end{enumerate}
    \vspace*{6cm}

    \item[\textbf{3.}] Un niño de \(65.0\,\text{kg}\) y su hermana de \(40.0\,\text{kg}\), ambos con patines, están frente a frente en reposo. La niña empuja duro al niño y lo envía hacia atrás con una velocidad de \(2.90\,\text{m/s}\) hacia el oeste. Ignore la fricción.
    \begin{enumerate}
        \item Describa el movimiento posterior de la niña.
        \item ¿Cuánta energía potencial del cuerpo de la niña se convierte en energía mecánica del sistema niño-niña?
        \item ¿Se conserva la cantidad de movimiento del sistema niño-niña en el proceso del empuje? Si es así, explique cómo ello es posible.
        \item Explique cómo es posible considerando que hay fuerzas grandes en acción.
        \item ¿Por qué no existe movimiento anticipado y mucho movimiento posterior?
    \end{enumerate}
    \vspace*{7cm}
\end{itemize}

\newpage

\begin{itemize}
    \item[\textbf{4.}] Dos bloques de masas \(m\) y \(3m\) se colocan sobre una superficie horizontal sin fricción. Un resorte ligero se une al bloque más masivo, y los bloques se juntan con el resorte entre ellos. Una cuerda que inicialmente mantiene a los bloques juntos se quema; después de esto, el bloque de masa \(3m\) se mueve hacia la derecha con una rapidez de \(2.00\,\text{m/s}\).
    \begin{enumerate}
        \item ¿Cuál es la velocidad del bloque de masa \(m\)?
        \item Encuentre la energía potencial elástica original del sistema, empleando \(m = 0.350\,\text{kg}\).
        \item ¿La energía original está en el resorte o en la cuerda?
        \item Explique su respuesta al inciso (c).
        \item ¿La cantidad de movimiento del sistema se conserva en el proceso de rompimiento?
        \item Explique cómo es posible considerando que hay fuerzas grandes en acción.
        \item ¿Por qué no existe movimiento anticipado y mucho movimiento posterior?
    \end{enumerate}
    \vspace*{8cm}

    \item[\textbf{5.}] Cuando usted salta recto hacia arriba tan alto como puede, ¿cuál es el orden de magnitud de la máxima rapidez de retroceso que le da a la Tierra? Modele la Tierra como un objeto perfectamente sólido. En su solución, establezca las cantidades físicas que toma como datos y los valores que mide o estima para ellos.
    \vspace*{6cm}
\end{itemize}

\subsection*{SECCIÓN 9.3 Análisis de modelos: Sistema no aislado (cantidad de movimiento)}

\begin{itemize}
    \item[\textbf{6.}] Una pelota de beisbol se aproxima a la placa de home con una rapidez de \(45.0\,\text{m/s}\), moviéndose horizontalmente justo antes de ser golpeada por un bat. El jugador batea la pelota dándole una rapidez de \(55.0\,\text{m/s}\). La bola tiene una masa de \(145\,\text{g}\) y está en contacto con el bat durante \(2.00\,\text{ms}\). ¿Cuál es el vector fuerza promedio que la bola ejerce sobre el bat durante su interacción?
    \vspace*{6cm}
\end{itemize}

\newpage

\begin{itemize}
    \item[\textbf{7.}] Un deslizador de masa \(m\) es libre de deslizarse a lo largo de un pista de aire horizontal. Se empuja contra un lanzador en un extremo de la pista. Modele el lanzador como un resorte ligero con constante elástica \(k\) comprimido una distancia \(x\). El deslizador se libera desde el reposo.
    \begin{enumerate}
        \item Muestre que el deslizador logra una rapidez de
        \[
        v = x\left(\frac{k}{m}\right)^{1/2}.
        \]
        \item Pruebe que la magnitud del impulso impartido al deslizador está dado por
        \[
        I = x(km)^{1/2}.
        \]
        \item ¿Se realiza más trabajo sobre un carrito con masa más grande o una más pequeña?
    \end{enumerate}
    \vspace*{6cm}

    \item[\textbf{8.}] Usted y su hermano discuten a menudo acerca de cómo sostener con seguridad a un niño en un auto en movimiento. Usted insiste en que los asientos especiales para niños son importantes en aumentar las posibilidades de que un niño sobreviva un accidente. Su hermano dice que, siempre y cuando su esposa se siente al lado de él con un cinturón de seguridad mientras conduce, puede sostener a su niño en su regazo en un accidente. Usted decide realizar un cálculo para tratar de convencer a su hermano. Considere una colisión hipotética en la que el niño de \(12\,\text{kg}\) y sus padres van en un auto que viaja a \(60\,\text{mi/h}\) en relación con el suelo. El auto golpea una pared, un árbol, u otro auto, y se detiene en \(0.10\,\text{s}\). Usted desea mostrarle a su hermano la magnitud de la fuerza necesaria para que su esposa sostenga a su hijo durante la colisión.
    \vspace*{6cm}

    \item[\textbf{9.}] El frente de \(1.20\,\text{m}\) de un auto de \(1400\,\text{kg}\) está diseñado como una “zona de deformación” que colapsa para absorber el impacto de una colisión. Si un auto que viaja a \(25.0\,\text{m/s}\) se detiene uniformemente en \(1.20\,\text{m}\),
    \begin{enumerate}
        \item ¿cuánto dura la colisión?
        \item ¿cuál es la magnitud de la fuerza promedio en relación con el auto?
        \item ¿cuál es la aceleración del auto? Exprese la aceleración como un múltiplo de la aceleración debida a la gravedad.
    \end{enumerate}
    \vspace*{6cm}

    \item[\textbf{10.}] La magnitud de la fuerza neta que se ejerce en la dirección \(x\) sobre una partícula de \(2.50\,\text{kg}\) varía en el tiempo como se muestra en la figura P9.10. Encuentre:
    \begin{enumerate}
        \item el impulso de la fuerza en el intervalo de tiempo entre \(0\) y \(5.00\,\text{s}\),
        \item la velocidad final de la partícula si originalmente está en reposo,
        \item la velocidad final si su velocidad original es \(-2.00\,\hat{i}\,\text{m/s}\),
        \item la fuerza promedio ejercida en la partícula durante el intervalo de tiempo entre \(0\) y \(5.00\,\text{s}\).
    \end{enumerate}
    \vspace*{6cm}
\end{itemize}

\newpage

\begin{itemize}
    \item[\textbf{11.}] Cae agua sin salpicar con una rapidez de \(0.250\,\text{L/s}\) desde una altura de \(2.60\,\text{m}\) en una cubeta de \(0.750\,\text{kg}\) sobre una báscula. Si la cubeta originalmente está vacía, ¿qué lee la báscula \(3.00\,\text{s}\) después de que el agua comienza a acumularse en ella?
    \vspace*{6cm}
\end{itemize}

\subsection*{SECCIÓN 9.4 Colisiones en una dimensión}

\begin{itemize}
    \item[\textbf{12.}] Un auto de \(1200\,\text{kg}\) viaja inicialmente a \(25.0\,\text{m/s}\) hacia el este y choca en la parte trasera de un camión de \(9000\,\text{kg}\) que se mueve en la misma dirección a \(20.0\,\text{m/s}\). La velocidad del auto inmediatamente después de la colisión es \(18.0\,\text{m/s}\) hacia el este.
    \begin{enumerate}
        \item ¿Cuál es la velocidad del camión inmediatamente después de la colisión?
        \item ¿Cuál es el cambio en la energía mecánica del sistema auto-camión en la colisión?
        \item Explique este cambio en energía mecánica.
    \end{enumerate}
    \vspace*{6cm}

    \item[\textbf{13.}] Un vagón de ferrocarril de \(2.50 \times 10^4\,\text{kg}\) de masa se mueve con una rapidez de \(4.00\,\text{m/s}\). Choca y se une con otros tres vagones acoplados, cada uno de la misma masa que el vagón solo y se mueven en la misma dirección con una rapidez inicial de \(2.00\,\text{m/s}\).
    \begin{enumerate}
        \item ¿Cuál es la rapidez de los cuatro vagones después de la colisión?
        \item ¿Cuánta energía mecánica se pierde en la colisión?
    \end{enumerate}
    \vspace*{6cm}

    \item[\textbf{14.}] Cuatro vagones, cada uno de \(2.50 \times 10^4\,\text{kg}\) de masa, se acoplan y avanzan a lo largo de pistas horizontales con rapidez \(v_i\) hacia el sur. Un actor de cine muy fuerte, que viaja en el segundo vagón, desacopla el vagón frontal y le da un gran empujón, lo que aumenta su rapidez a \(4.00\,\text{m/s}\) hacia el sur. Los tres vagones restantes continúan moviéndose hacia el sur, ahora a \(2.00\,\text{m/s}\).
    \begin{enumerate}
        \item Encuentre la rapidez inicial de los cuatro vagones.
        \item ¿Qué tanto cambió la energía potencial química en el cuerpo del actor?
        \item Establezca la relación entre el proceso aquí descrito y el proceso del problema 13.
    \end{enumerate}
    \vspace*{6cm}
\end{itemize}

\newpage

\begin{itemize}
    \item[\textbf{15.}] Un automóvil de masa \(m_1\) que se mueve con rapidez \(v_1\) colisiona y se acopla con la parte trasera de un camión de masa \(2m_1\) que inicialmente viaja en la misma dirección que el automóvil, pero a una menor rapidez \(v_2\).
    \begin{enumerate}
        \item ¿Cuál es la rapidez \(v_f\) de los dos vehículos inmediatamente después de la colisión?
        \item ¿Cuál es el cambio de energía cinética del sistema automóvil-camión en la colisión?
    \end{enumerate}
    \vspace*{6cm}

    \item[\textbf{16.}] Una bala de \(7.00\,\text{g}\), cuando se dispara desde un arma en un bloque de madera de \(1.00\,\text{kg}\) sostenido en un tornillo de banco, penetra el bloque a una profundidad de \(8.00\,\text{cm}\). Este bloque de madera se coloca sobre una superficie horizontal sin fricción, y una segunda bala de \(7.00\,\text{g}\) se dispara desde el arma en el bloque. En este caso, ¿a qué profundidad penetra la bala en el bloque?
    \vspace*{6cm}

    \item[\textbf{17.}] Una pelota de tenis de \(57.0\,\text{g}\) de masa se sostiene justo arriba de un balón de basquetbol de \(590\,\text{g}\) de masa. Con sus centros verticalmente alineados, ambos se liberan desde el reposo en el mismo momento, para caer una distancia de \(1.20\,\text{m}\), como se muestra en la figura P9.17.
    \begin{enumerate}
        \item Encuentre la magnitud de la velocidad hacia abajo con la que el balón llega al suelo.
        \item Suponga que una colisión elástica con el suelo instantáneamente invierte la velocidad del balón mientras la pelota de tenis aún se mueve hacia abajo. Luego las dos bolas se encuentran en una colisión elástica. ¿A qué altura rebota la pelota de tenis?
    \end{enumerate}
    \vspace*{6cm}

    \item[\textbf{18.}] Tres carros de masas
    \[
    m_1 = 4.00\,\text{kg}, \quad
    m_2 = 10.0\,\text{kg}, \quad
    m_3 = 3.00\,\text{kg}
    \]
    se mueven sobre una pista horizontal sin fricción con magnitudes de velocidad
    \[
    v_1 = 5.00\,\text{m/s} \text{ hacia la derecha},
    \]
    \[
    v_2 = 3.00\,\text{m/s} \text{ hacia la derecha},
    \]
    \[
    v_3 = 4.00\,\text{m/s} \text{ hacia la izquierda},
    \]
    como se muestra en la figura P9.18. Acopladores de velcro hacen que los carros queden unidos después de chocar.
    \begin{enumerate}
        \item Encuentre la velocidad final del tren unido de tres carros.
        \item ¿Qué pasaría si? ¿Su respuesta requiere que todos los carros choquen y se unan en el mismo momento? ¿Qué sucede si chocan en diferente orden?
    \end{enumerate}
    \vspace*{6cm}
\end{itemize}

\newpage

\subsection*{SECCIÓN 9.5 Colisiones en dos dimensiones}

\begin{itemize}
    \item[\textbf{19.}] Usted ha sido contratado como un testigo experto por un abogado para un juicio que involucró un accidente de tráfico. El cliente del abogado, el demandante en este caso, viajaba al este hacia una intersección a \(13.0\,\text{m/s}\) como se midió justo antes del accidente por un medidor de velocidad de carretera, el cual era un testigo digno de confianza. Cuando el demandante entró en la intersección, su auto fue golpeado por un conductor en dirección norte, el acusado en este caso conducía un auto con masa idéntica a la del demandante. Los vehículos quedaron pegados juntos después de la colisión y la izquierda paralela al norte del este, medido por investigadores de accidentes. El acusado afirma que él viajaba dentro del límite de velocidad de \(35\,\text{mi/h}\). ¿Qué le aconseja usted al abogado?
    \vspace*{6cm}

    \item[\textbf{20.}] Dos discos de juego, de igual masa, uno anaranjado y el otro amarillo, están involucrados en una colisión oblicua elástica. El disco amarillo inicialmente está en reposo y es golpeado por el disco anaranjado que se mueve con una rapidez de \(5.00\,\text{m/s}\). Después de la colisión, el disco anaranjado se mueve a lo largo de una dirección que forma un ángulo de \(37.0^\circ\) con su dirección de movimiento inicial. Las velocidades de los dos discos son perpendiculares después de la colisión. Determine la rapidez final de cada disco.
    \vspace*{6cm}

    \item[\textbf{21.}] Dos discos de juego, de igual masa, uno anaranjado y el otro amarillo, están involucrados en una colisión oblicua elástica. El disco amarillo inicialmente está en reposo y es golpeado por el disco anaranjado que se mueve con rapidez \(v_0\). Después de la colisión, el disco anaranjado se mueve a lo largo de una dirección que forma un ángulo \(\theta\) con su dirección de movimiento inicial. Las velocidades de los dos discos son perpendiculares después de la colisión. Determine la rapidez final de cada disco.
    \vspace*{6cm}

    \item[\textbf{22.}] Un corredor de futbol americano, de \(90.0\,\text{kg}\), se mueve al este con una rapidez de \(5.00\,\text{m/s}\) y es detenido por un oponente de \(95.0\,\text{kg}\) que corre al norte con una rapidez de \(3.00\,\text{m/s}\).
    \begin{enumerate}
        \item Explique por qué la tacleada exitosa constituye una colisión perfectamente inelástica.
        \item Calcule la velocidad de los jugadores inmediatamente después de la tacleada.
        \item Determine la energía mecánica que se disipa como resultado de la colisión. Explique la energía perdida.
    \end{enumerate}
    \vspace*{6cm}
\end{itemize}

\newpage

\begin{itemize}
    \item[\textbf{23.}] Un protón, que se mueve con una velocidad de \(v_i\), colisiona elásticamente con otro protón que inicialmente está en reposo. Suponiendo que los protones tienen la misma rapidez después de la colisión, encuentre:
    \begin{enumerate}
        \item la rapidez de cada protón después de la colisión en términos de \(v_i\),
        \item la dirección de los vectores velocidad después de la colisión.
    \end{enumerate}
    \vspace*{6cm}
\end{itemize}

\subsection*{SECCIÓN 9.6 El centro de masa}

\begin{itemize}
    \item[\textbf{24.}] A una pieza uniforme de hoja metálica se le da la forma que se muestra en la figura P9.24. Calcule las coordenadas \(x\) y \(y\) del centro de masa de la pieza.
    \vspace*{6cm}

    \item[\textbf{25.}] En la selva, los exploradores encuentran un antiguo monumento en forma de triángulo isósceles, como se muestra en la figura P9.25. El monumento está hecho de miles de pequeños bloques de piedra de densidad \(3800\,\text{kg/m}^3\). La altura del monumento es \(15.7\,\text{m}\) y su base \(64.8\,\text{m}\) en su base y grosor de \(3.60\,\text{m}\) desde el frente hasta la parte trasera. Antes de construir el monumento, todos los bloques de piedra yacían en el suelo. ¿Cuánto trabajo hicieron los constructores sobre los bloques para colocarlos en posición durante la edificación del monumento?
    \vspace*{6cm}

    \item[\textbf{26.}] Una barra de \(30.0\,\text{cm}\) de longitud tiene densidad lineal (masa por longitud) dada por
    \[
    \lambda = 50.0 + 20.0x
    \]
    donde \(x\) es la distancia desde un extremo, medida en metros, y \(\lambda\) está en gramos/metro.
    \begin{enumerate}
        \item ¿Cuál es la masa de la barra?
        \item ¿A qué distancia del extremo \(x=0\) está su centro de masa?
    \end{enumerate}
    \vspace*{6cm}
\end{itemize}

\newpage

\subsection*{SECCIÓN 9.7 Movimiento de un sistema de partículas}

\begin{itemize}
    \item[\textbf{27.}] Considere un sistema de dos partículas en el plano \(xy\):
    \[
    m_1 = 2.00\,\text{kg}
    \]
    está en la ubicación
    \[
    \vec r_1 = (1.00\hat{i} + 2.00\hat{j})\,\text{m}
    \]
    y tiene una velocidad de
    \[
    (3.00\hat{i} + 0.500\hat{j})\,\text{m/s};
    \]
    \[
    m_2 = 3.00\,\text{kg}
    \]
    está en
    \[
    \vec r_2 = (-4.00\hat{i} - 3.00\hat{j})\,\text{m}
    \]
    y tiene velocidad
    \[
    (3.00\hat{i} - 2.00\hat{j})\,\text{m/s}.
    \]
    \begin{enumerate}
        \item Grafique esas partículas en una hoja cuadriculada usando sus vectores de posición y muestre la posición del centro de masa del sistema.
        \item Determine la velocidad del centro de masa.
        \item Determine la cantidad de movimiento lineal total del sistema.
    \end{enumerate}
    \vspace*{6cm}

    \item[\textbf{28.}] El vector de posición de una partícula de \(3.50\,\text{g}\) que se mueve en el plano \(xy\) varía en el tiempo de acuerdo con
    \[
    \vec r_1 = (3\hat{i} + 3t\hat{j} + 2t^2\hat{k})
    \]
    donde \(t\) está en segundos y \(\vec r\) está en centímetros. Al mismo tiempo, el vector de posición de una partícula de \(5.50\,\text{g}\) varía como
    \[
    \vec r_2 = 3\hat{i} - 2t\hat{j} - 6t\hat{k}.
    \]
    En \(t = 2.50\,\text{s}\), determine:
    \begin{enumerate}
        \item el vector de posición del centro de masa,
        \item la cantidad de movimiento lineal,
        \item la velocidad del centro de masa,
        \item la aceleración del centro de masa,
        \item la fuerza neta que se ejerce en el sistema de dos partículas.
    \end{enumerate}
    \vspace*{7cm}

    \item[\textbf{29.}] Usted ha sido contratado como testigo experto en una investigación de un incidente de un dron con cuatro soportes. El incidente ocurrió durante una lluvia de meteoritos muy rara durante la cual varios trozos masivos de metal meteórico pasaron a través de la atmósfera y golpearon el suelo. El dron no tripulado estaba rondando en reposo en el centro de una casa en llamas y acababa de caer lentamente de fuego, cuando explotó espontáneamente en cuatro grandes trozos. Las ubicaciones de las cuatro piezas en el suelo se miden como sigue, en relación con el centro de la casa sobre la que el dron estaba rondando:
    \[
    \begin{array}{c|c|c|c}
    \text{Pieza} & \text{Masa (kg)} & \text{Distancia desde el centro de la casa (m)} & \text{Dirección desde la casa} \\
    \hline
    1 & 80.0 & 150 & \text{Hacia el oeste} \\
    2 & 120 & 75.0 & \text{Hacia el norte} \\
    3 & 50.0 & 90.0 & 20.0^\circ \text{ al oeste del sur} \\
    4 & 150 & 50.0 & 20.0^\circ \text{ al norte del este}
    \end{array}
    \]
    El Departamento de Bomberos está sugiriendo que el dron estaba defectuoso y explotó mientras estaba en uso. El fabricante sugiere que el dron fue golpeado por un meteorito, causando la explosión. Realice un cálculo que muestre la evidencia que sugiera un acuerdo con una de estas posiciones.
    \vspace*{7cm}
\end{itemize}

\newpage

\begin{itemize}
    \item[\textbf{30.}] Para un proyecto de tecnología, un estudiante construyó un vehículo de \(6.00\,\text{kg}\) de masa total, que se mueve por sí solo. Como se muestra en la figura P9.30, corre sobre dos orugas ligeras que pasan alrededor de cuatro ruedas ligeras. Un carrete se acopla a uno de los ejes, y una cuerda originalmente enrollada sobre el carrete pasa sobre una polea unida al vehículo para soportar una carga elevada. Después de que el vehículo se libera desde el reposo, la carga desciende lentamente, desenrolla la cuerda para girar el eje y hace que el vehículo se mueva hacia adelante. El carrete tiene una forma cónica de modo que la carga desciende a una rapidez baja constante mientras el vehículo se mueve horizontalmente a través del suelo con aceleración constante, y alcanza una velocidad final de \(3.00\,\hat{i}\,\text{m/s}\).
    \begin{enumerate}
        \item ¿El suelo ejerce un impulso sobre el vehículo? Si es así, ¿cuánto es?
        \item ¿El suelo realiza trabajo sobre el vehículo?
        \item ¿Cuál es la procedencia de la cantidad de movimiento del vehículo?
        \item ¿Cuál es la procedencia de la energía cinética final del vehículo?
        \item Identifique la fuerza específica que causa la aceleración hacia adelante del vehículo.
    \end{enumerate}
    \vspace*{9cm}
\end{itemize}

\end{document}
